% Generated draft: requires mathematical and source review.
\documentclass[11pt]{article}
\usepackage[margin=25mm]{geometry}
\usepackage{fontspec}
\setmainfont{Latin Modern Roman}
\usepackage{amsmath,amssymb,mathtools,unicode-math}
\setmathfont{Latin Modern Math}
\usepackage{enumitem,needspace,xurl,hyperref}
\setlength{\parindent}{0pt}
\setlength{\parskip}{5pt}
\setlength{\emergencystretch}{3em}
\newcommand{\N}{\mathbb{N}}
\newcommand{\Z}{\mathbb{Z}}
\newcommand{\Q}{\mathbb{Q}}
\newcommand{\R}{\mathbb{R}}
\newcommand{\C}{\mathbb{C}}
\newcommand{\F}{\mathbb{F}}
\newcommand{\A}{\mathbb{A}}
\newcommand{\PP}{\mathbb{P}}
\newcommand{\E}{\mathbb{E}}
\newcommand{\eps}{\varepsilon}
\newcommand{\dd}{\,\mathrm{d}}
\begin{document}
\section*{UnsolvedMath ID: 30006368}
\textbf{Source problem code:} catalogue entry
\subsection{Definitions and mathematical statement}
A closed manifold means a compact smooth manifold without boundary. The supplied record asks about fixed-energy magnetic geodesics on a finite-dimensional closed manifold, but does not specify the magnetic data or the hypotheses needed for its intended connectivity assertion.

For a conventional formulation, suppose that a magnetic system consists of a Riemannian manifold \((M,g)\) and a closed real differential \(2\)-form \(\sigma\). This is an explicit interpretation of the terminology, rather than additional hypotheses established by the supplied passage. Define the Lorentz-force endomorphism \(Y:TM\to TM\) by
\[
 g_x(Y_xv,w)=\sigma_x(v,w),\qquad x\in M,\quad v,w\in T_xM.
\]
Writing \(\nabla\) for the Levi--Civita connection of \(g\), a magnetic geodesic in this convention is a \(C^2\) curve \(\gamma:[0,T]\to M\) satisfying
\[
 \nabla_{\dot\gamma}\dot\gamma=Y_{\gamma}\dot\gamma.
\]
Its kinetic energy is \(E(\gamma,\dot\gamma)=\tfrac12 g_{\gamma}(\dot\gamma,\dot\gamma)\).

The extracted question asks under which hypotheses on the magnetic system and energy the following endpoint-connectivity property holds:
\[
 \begin{gathered}
 \forall p,q\in M\quad \forall k>0\quad
 \exists T>0\quad \exists\gamma\in C^2([0,T],M):\\
 \gamma(0)=p,\qquad \gamma(T)=q,\\
 \nabla_{\dot\gamma(t)}\dot\gamma(t)=Y_{\gamma(t)}\dot\gamma(t),\qquad
 \tfrac12 g_{\gamma(t)}(\dot\gamma(t),\dot\gamma(t))=k
 \quad(0\le t\le T).
 \end{gathered}
\]
Here the travel time is allowed to vary. The displayed formula expresses the unrestricted positive-energy question; it is not asserted to hold. The source also asks how the admissible energy range depends on Ma\~n\'e critical values. Its supplied dated assessment singles out a supercritical threshold question, but neither identifies the critical value nor specifies a corresponding inequality or the hypotheses under which that value is defined. Consequently, a precise supercritical specialization cannot be recovered from this record.

Further qualifications remain unresolved. Connectedness is not stated: endpoints in distinct connected components cannot be joined by a curve, so any intended universal assertion needs a connectedness hypothesis or a restriction on endpoints. The passage also does not say whether the two endpoints must be distinct; including \(p=q\) requires a positive-energy return to the same point. These choices must be settled before the displayed quantifiers can serve as the definitive problem statement.
\subsection{Short English statement}
For a magnetic system on a closed manifold, ask when every prescribed pair of endpoints can be joined by a magnetic geodesic with prescribed positive kinetic energy and freely chosen travel time. Determine how this property depends on the relevant Ma\~n\'e critical values. The magnetic hypotheses, endpoint restrictions, and intended critical threshold are missing, so the formulation remains provisional. The supplied open status is dated curation and requires independent review.
\subsection{Sources}
\begin{itemize}
\item Local registry: \texttt{data/unsolvedmath/problems.json}, supplied record for permanent UnsolvedMath ID \(30006368\), exact source problem code \texttt{catalogue entry}, titled ``Two-Point Connectivity by Fixed-Energy Magnetic Geodesics.'' No individual proposer is supplied; the proposed year is recorded as 2025.
\item Supplied report citation: \emph{Dynamische Systeme (2026)}, Oberwolfach Reports, DOI 10.4171/owr/2025/31, pp.\ 1647--1706. \url{https://doi.org/10.4171/owr/2025/31}. The record attributes the extracted question to this workshop-report context.
\item Supplied literature source, labeled ``Albers--Maier primary paper'': \url{https://doi.org/10.1007/s00220-026-05651-z}. No title or precise supporting result is supplied.
\item Supplied literature source, labeled ``Oberwolfach report'': \url{https://doi.org/10.4171/owr/2025/31}.
\item Literature assessment dated exactly \texttt{2026-08-22}: the supplied curation labels the finite-dimensional supercritical Ma\~n\'e-threshold connectivity question on closed manifolds open and says the extracted passage lacks hypotheses sufficient for a stronger classification. This is supplied dated curation, not an independently verified assessment.
\item The original passage refers to reference [14] for Ma\~n\'e critical values; its bibliographic identity is not supplied. The registry labels the statement \texttt{corrected\_verified} and reports that the official report was checked and notation repaired. These are supplied verification claims, not verification performed for this draft.
\end{itemize}
\end{document}
